\section{Modelado del sistema en ArchiMate}\label{sec:archimate-modelado}

ArchiMate es un lenguaje de modelado estandarizado por~\textit{The Open Group} que permite representar de manera estructurada los componentes y relaciones que conforman una arquitectura empresarial. Su utilización permite establecer una visión integrada de los elementos de negocio, aplicación y tecnología, facilitando la representación de las relaciones existentes entre las necesidades de la organización, los servicios de información y la infraestructura que los soporta~\citep{opengroup_archimate_overview}.

En el presente proyecto, ArchiMate se utilizó para modelar la arquitectura propuesta para el servicio de directorio del grupo de investigación \ac{GRID}. El modelado permitió representar la relación entre los actores que hacen uso de los servicios, los servicios tecnológicos que requieren mecanismos de autenticación y autorización, y los componentes que conforman la solución de directorio seleccionada. De esta manera, se estableció una representación de la solución que permite visualizar cómo el servicio de directorio se integra con los diferentes servicios considerados dentro del alcance del proyecto.

El modelo contempla principalmente los elementos relacionados con la gestión centralizada de identidades y el acceso a los servicios integrados, así como los componentes tecnológicos necesarios para proporcionar dicha funcionalidad. En este contexto, se representa a FreeIPA como componente central de la solución, encargado de proporcionar los servicios de directorio y gestión de identidades sobre los cuales se soporta la autenticación de los servicios integrados.

La representación mediante ArchiMate permitió, además, establecer las relaciones entre los diferentes elementos de la arquitectura y proporcionar una visión general de la solución, facilitando la comprensión de su estructura y de las interacciones entre sus componentes. En los Anexos~\ref{tab:elementos-negocio-archimate}, \ref{tab:elementos-aplicacion-archimate}, \ref{tab:elementos-tecnologicos-archimate}, \ref{tab:elementos-estrategia-archimate} y \ref{tab:relaciones-archimate} se especifica el significado de los principales elementos y relaciones utilizados en el modelado.
